\section{Experimental Setup}
\label{sec:experiments}

% OWNER: Jackson Wilke (methods/setup per proposal)

\subsection{Training}
\label{sec:training}
Both architectures are trained from scratch on each noise condition. The
vocabulary (10K types) is built once from the clean corpus and shared across
every condition and both models, so vocabulary never confounds the comparison.

The LSTM uses embedding and hidden dimension 256, two layers, dropout 0.2, and
tied input/output embeddings, optimized with Adam (learning rate $10^{-3}$,
batch size 64, gradient clipping at 1.0) for 5 epochs over 150K sentences per
condition. The RNNG (top-down, fixed-stack) uses word and hidden dimension 256,
two layers, dropout 0.3, Adam (learning rate $10^{-3}$, batch size 128 with a
15K-token cap), for 8 epochs over the same 150K sentences; per-verb surprisals
are obtained by word-synchronous beam search (beam size 100, word-beam 20).
Checkpoints are named by condition and seed. A full hyperparameter table is given
in Appendix~\ref{app:repro}.

\subsection{Conditions and Runs}
\label{sec:runs}
The grid is 5 noise rates $\times$ 2 architectures. The LSTM sweep was run on
two machines (local Apple MPS and a Colab T4) across three seeds, giving six
runs per condition; reported LSTM values aggregate over these runs, with $\pm$
denoting standard deviation. The RNNG arm contributes one full-scale seed of
record per condition, matched to the LSTM protocol; additional RNNG seeds are
pending, so we report RNNG point estimates without error bars.

\subsection{Trajectory Analysis}
\label{sec:trajectory}
To measure \emph{when} a rule is acquired, not only whether, we checkpoint the
LSTM every 500 optimizer steps and re-score the full minimal-pair set at each
checkpoint, yielding accuracy and attractor-gap trajectories over training
(Section~\ref{sec:results}).

\subsection{Hypotheses}
\label{sec:hypotheses}
We test three preregistered hypotheses:
\begin{itemize}
  \item \textbf{H1.} Overall minimal-pair accuracy decreases as the training
    noise rate increases.
  \item \textbf{H2.} The attractor gap $\Delta$ widens with noise as models lean
    increasingly on the linear nearest-noun cue.
  \item \textbf{H3.} The RNNG's attractor gap widens \emph{more slowly} than the
    LSTM's; an explicit structural bias buys robustness to noisy supervision.
\end{itemize}
We assess H3 by comparing the LSTM and RNNG attractor gaps across conditions on
their shared held-out set (Section~\ref{sec:metrics}). The LSTM gap carries a
seed-level mean $\pm$ std over its six runs; the RNNG arm is a single full-scale
seed, so we read the per-condition gap difference
$\Delta_{\text{LSTM}} - \Delta_{\text{RNNG}}$ as a direction to be tightened by
the pending multi-seed RNNG replication rather than as a formal test.
